%% CE 49X - Introduction to Computational Thinking and Data Science
%% Midterm Exam - Part 1 (Paper-Based)
%% Weeks 1--7
%% Instructor: Dr. Eyuphan Koc
%% Bogazici University - Spring 2026

\documentclass[11pt,a4paper]{article}

%% ==================== PACKAGES ====================
\usepackage[utf8]{inputenc}
\usepackage[T1]{fontenc}
\usepackage{amsmath,amssymb,amsfonts}
\usepackage{graphicx}
\usepackage[margin=0.7in,headheight=14pt]{geometry}
\usepackage{booktabs}
\usepackage{array}
\usepackage{enumerate}
\usepackage{xcolor}
\usepackage{fancyhdr}
\usepackage{listings}
\usepackage{courier}

%% Define custom colors
\definecolor{bunavy}{RGB}{0,20,60}
\definecolor{softblue}{RGB}{100,140,180}
\definecolor{codebg}{RGB}{245,245,245}

%% Listings style for Python
\lstset{
  language=Python,
  basicstyle=\ttfamily\small,
  backgroundcolor=\color{codebg},
  frame=single,
  framerule=0.5pt,
  rulecolor=\color{gray!50},
  keywordstyle=\color{blue!70!black},
  stringstyle=\color{red!60!black},
  commentstyle=\color{green!50!black},
  showstringspaces=false,
  breaklines=true,
  numbers=none,
  xleftmargin=4pt,
  xrightmargin=4pt,
  aboveskip=8pt,
  belowskip=8pt
}

%% Set up header and footer
\pagestyle{fancy}
\fancyhf{}
\lhead{CE 49X - Spring 2026}
\chead{Midterm Exam -- Part 1}
\rhead{Dr. E. Koc}
\lfoot{}
\cfoot{\thepage}
\rfoot{}

%% ==================== DOCUMENT BEGINS ====================
\begin{document}

%% Title and Instructions
\begin{center}
\Large\textbf{CE 49X: Introduction to Computational Thinking and Data Science}\\[0.1cm]
\large\textbf{Midterm Exam -- Part 1 (Paper-Based)}\\[0.05cm]
\normalsize Weeks 1--7 Material\\[0.2cm]
\end{center}

\noindent
\begin{tabular}{ll}
\textbf{Name:} \rule{8cm}{0.4pt} & \textbf{Student ID:} \rule{4cm}{0.4pt}\\[0.2cm]
\textbf{Date:} April 8, 2026 & \textbf{Duration:} 60 minutes
\end{tabular}

\vspace{0.2cm}
\noindent\rule{\linewidth}{1pt}

\noindent\textbf{Instructions:}
\begin{itemize}
\setlength\itemsep{0em}
\item This is Part 1 of the midterm exam (paper-based).
\item Answer all questions. Show your work for full credit on written questions.
\item One page (both sides) of handwritten notes is allowed.
\item No calculators or electronic devices.
\item A helpful formulas page is provided at the end.
\item Total points: \textbf{50}
\end{itemize}

\noindent\rule{\linewidth}{1pt}
\vspace{0.2cm}

%% ==================== SECTION A: MULTIPLE CHOICE ====================
\section*{Section A: Multiple Choice \hfill [20 points]}

\noindent \textit{Circle the single best answer for each question. Each question is worth 2 points.}

\vspace{0.3cm}

\noindent\textbf{1.} What does the following expression evaluate to?
\begin{lstlisting}
len([x for x in range(10) if x % 3 == 0])
\end{lstlisting}
\begin{enumerate}[(a)]
\item 3
\item 4
\item 10
\item 0
\end{enumerate}

\vspace{0.3cm}

\noindent\textbf{2.} Which import statement allows you to call \lstinline{np.array([1, 2, 3])}?
\begin{enumerate}[(a)]
\item \lstinline{import numpy}
\item \lstinline{from numpy import array}
\item \lstinline{import numpy as np}
\item \lstinline{import np from numpy}
\end{enumerate}

\vspace{0.3cm}

\noindent\textbf{3.} Given \lstinline{a = np.array([1, 2, 3, 4, 5])}, what does \lstinline{a[1:4]} return?
\begin{enumerate}[(a)]
\item \lstinline{array([1, 2, 3, 4])}
\item \lstinline{array([2, 3, 4])}
\item \lstinline{array([2, 3, 4, 5])}
\item \lstinline{array([1, 2, 3])}
\end{enumerate}

\vspace{0.3cm}

\noindent\textbf{4.} Which pandas method fills missing values with the \textbf{previous} row's value?
\begin{enumerate}[(a)]
\item \lstinline{df.dropna()}
\item \lstinline{df.fillna(method='ffill')}
\item \lstinline{df.interpolate()}
\item \lstinline{df.fillna(0)}
\end{enumerate}

\vspace{0.3cm}

\noindent\textbf{5.} Which visual encoding channel is \textbf{most accurate} for comparing quantitative values?
\begin{enumerate}[(a)]
\item Color intensity
\item Angle
\item Area
\item Position
\end{enumerate}

\newpage

\noindent\textbf{6.} According to Tufte's principles, what should a well-designed chart maximize?
\begin{enumerate}[(a)]
\item Number of colors used
\item Data-ink ratio
\item Chart area
\item Number of gridlines
\end{enumerate}

\vspace{0.3cm}

\noindent\textbf{7.} A dataset has mean $\mu = 50$ and standard deviation $\sigma = 8$. What is the z-score of a value of 66?
\begin{enumerate}[(a)]
\item 1.5
\item 2.0
\item 2.5
\item 8.0
\end{enumerate}

\vspace{0.3cm}

\noindent\textbf{8.} A hypothesis test yields a p-value of 0.03 at significance level $\alpha = 0.05$. What is the correct conclusion?
\begin{enumerate}[(a)]
\item There is a 3\% chance that $H_0$ is true
\item Reject $H_0$
\item Fail to reject $H_0$
\item The effect size is 0.03
\end{enumerate}

\vspace{0.3cm}

\noindent\textbf{9.} A model predicts ``safe'' for every sample in a dataset where 95\% of samples are actually safe. Its accuracy is 95\%. Why is this misleading?
\begin{enumerate}[(a)]
\item The model is overfitting
\item It never identifies the dangerous 5\%
\item The training set is too small
\item Accuracy requires cross-validation
\end{enumerate}

\vspace{0.3cm}

\noindent\textbf{10.} In Naive Bayes, what does the ``naive'' assumption mean?
\begin{enumerate}[(a)]
\item The model uses no prior information
\item Features are conditionally independent given the class
\item The model always predicts the most common class
\item Training data must be normally distributed
\end{enumerate}

\newpage
%% ==================== SECTION B: WRITTEN QUESTIONS ====================
\section*{Section B: Written Questions \hfill [30 points]}

%% ==================== QUESTION 1 ====================
\subsection*{Question 1: Code Tracing \& Debugging \hfill [6 points]}

\noindent\textbf{(a) (3 points)} What is the \textbf{exact} output of the following code?

\begin{lstlisting}
beams = [3.5, 7.2, 4.8, 12.1, 6.3]
long_beams = [b for b in beams if b > 5.0]
total = sum(long_beams)
count = len(long_beams)
print(f"Found {count} beams totaling {total:.1f} m")
print(f"Average: {total/count:.2f} m")
\end{lstlisting}

\vspace{4.5cm}

\noindent\textbf{(b) (3 points)} The following function has a bug. Identify the error, explain what the code currently prints, and write the corrected line(s).

\begin{lstlisting}
def classify_soil(spt_n_values):
    """Classify as 'soft' (N<10), 'medium' (10<=N<30), 'hard' (N>=30)"""
    results = {}
    for n in spt_n_values:
        if n < 10:
            results[n] = 'soft'
        elif n < 30:
            results[n] = 'medium'
        else:
            results[n] = 'hard'
    return results

data = [5, 25, 5, 40, 25]
print(len(classify_soil(data)))
\end{lstlisting}

\vspace{5cm}

%% ==================== QUESTION 2 ====================
\subsection*{Question 2: NumPy \& Pandas \hfill [6 points]}

\noindent\textbf{(a) (3 points)} Explain the difference between a NumPy \textbf{view} and a \textbf{copy}. Given the code below, what is the value of \lstinline{speeds[1]} after execution? Explain why.

\begin{lstlisting}
import numpy as np
speeds = np.array([40, 55, 30, 65, 20])
subset = speeds[1:4]
subset[0] = 99
\end{lstlisting}

\vspace{4.5cm}

\noindent\textbf{(b) (3 points)} Write a \textbf{single} pandas expression that filters a DataFrame \lstinline{df} to rows where \lstinline{AVERAGE_SPEED < 20} AND \lstinline{NUMBER_OF_VEHICLES > 100}, and returns only the \lstinline{GEOHASH} and \lstinline{AVERAGE_SPEED} columns.

\vspace{4.5cm}

\newpage
%% ==================== QUESTION 3 ====================
\subsection*{Question 3: Visualization \hfill [6 points]}

\noindent\textbf{(a) (3 points)} A colleague prepares a chart for a municipal earthquake damage report. The chart is a \textbf{3D pie chart} using the \textbf{jet (rainbow) colormap} to show the percentage of damaged buildings across 8 districts. Identify \textbf{two} specific problems with this chart and suggest a fix for each.

\vspace{5.5cm}

\noindent\textbf{(b) (3 points)} What is \textbf{Anscombe's Quartet}? In 2--3 sentences, explain what it demonstrates and the lesson it teaches for data analysis.

\vspace{5cm}

%% ==================== QUESTION 4 ====================
\subsection*{Question 4: Statistical Analysis \hfill [6 points]}

\noindent\textbf{(a) (3 points)} A batch of 28-day concrete cylinder tests gives a mean compressive strength of $\mu = 32.5$ MPa with a standard deviation of $\sigma = 3.2$ MPa. Compute the z-score for a cylinder that tested at 26.1 MPa. Is this result ``extreme'' by the $|z| > 2$ convention?

\vspace{4.5cm}

\noindent\textbf{(b) (3 points)} Define \textbf{Type I error} and \textbf{Type II error}. For a concrete strength test where $H_0$: ``the concrete meets the 30 MPa specification,'' describe what each error type means in practical engineering terms. Which is more dangerous for structural safety?

\vspace{5.5cm}

\newpage
%% ==================== QUESTION 5 ====================
\subsection*{Question 5: Machine Learning \& Naive Bayes \hfill [6 points]}

\noindent\textbf{(a) (3 points)} A soil liquefaction classifier is evaluated on 200 test sites. The confusion matrix is:

\begin{center}
\begin{tabular}{l|cc}
 & \textbf{Predicted Stable} & \textbf{Predicted Liquefiable} \\
\hline
\textbf{Actually Stable} & 160 & 15 \\
\textbf{Actually Liquefiable} & 5 & 20 \\
\end{tabular}
\end{center}

Compute the \textbf{precision} and \textbf{recall} for the ``Liquefiable'' class. Should the engineer prioritize precision or recall in this application? Justify in 1--2 sentences.

\vspace{5.5cm}

\noindent\textbf{(b) (3 points)} A bridge inspector uses Bayes' theorem to assess beam condition. Given:
\begin{itemize}
\setlength\itemsep{0em}
\item Prior: $P(\text{Good}) = 0.70$, $P(\text{Poor}) = 0.30$
\item Likelihood of observing an 8mm crack: $P(\text{8mm} \,|\, \text{Good}) = 0.02$, $P(\text{8mm} \,|\, \text{Poor}) = 0.15$
\end{itemize}

\noindent Compute $P(\text{Poor} \,|\, \text{8mm crack})$. Show your work.

\vspace{6cm}

\newpage
%% ==================== HELPFUL FORMULAS ====================
\section*{Helpful Formulas}

\small

\textbf{Statistics:}
\begin{align*}
&\text{Z-score: } z = \frac{x - \mu}{\sigma} \\[0.3cm]
&\text{68-95-99.7 Rule: } P(\mu \pm 1\sigma) \approx 68\%, \quad P(\mu \pm 2\sigma) \approx 95\%, \quad P(\mu \pm 3\sigma) \approx 99.7\%
\end{align*}

\textbf{Classification Metrics:}
\begin{align*}
&\text{Accuracy} = \frac{TP + TN}{TP + TN + FP + FN} \\[0.2cm]
&\text{Precision} = \frac{TP}{TP + FP} \qquad \text{Recall} = \frac{TP}{TP + FN} \\[0.2cm]
&F_1 = \frac{2 \cdot \text{Precision} \cdot \text{Recall}}{\text{Precision} + \text{Recall}}
\end{align*}

\textbf{Confusion Matrix Layout:}
\begin{center}
\begin{tabular}{l|cc}
 & \textbf{Predicted Negative} & \textbf{Predicted Positive} \\
\hline
\textbf{Actually Negative} & TN & FP \\
\textbf{Actually Positive} & FN & TP \\
\end{tabular}
\end{center}

\textbf{Bayes' Theorem:}
$$P(A \,|\, B) = \frac{P(B \,|\, A) \cdot P(A)}{P(B)}$$
$$\text{where } P(B) = \sum_i P(B \,|\, A_i) \cdot P(A_i)$$

\textbf{Naive Bayes:}
$$P(L \,|\, f_1, \ldots, f_n) \propto P(L) \prod_{i=1}^{n} P(f_i \,|\, L)$$

\textbf{Gaussian PDF:}
$$P(f_i \,|\, L) = \frac{1}{\sqrt{2\pi\sigma_L^2}} \exp\!\left(-\frac{(f_i - \mu_L)^2}{2\sigma_L^2}\right)$$

\vspace{0.5cm}
\textbf{Visual Encoding Accuracy Hierarchy:}\\
Position $>$ Length $>$ Angle $>$ Area $>$ Color Intensity

\vspace{0.5cm}
\begin{center}
\noindent\rule{0.5\linewidth}{0.4pt}\\[0.3cm]
\textbf{End of Part 1}\\[0.2cm]
\textit{Please review your answers if time permits.}
\end{center}

\end{document}
